% Section 7.2, from lectures/L07/README.md: the goals, the evaluation questions, and the next
% lecture.

\needspace{12\baselineskip}
\section{Review}
\label{c7:sec:review}

\needspace{6\baselineskip}
\subsection*{What you should be able to do now}
After this chapter, you should be able to:
\begin{itemize}
  \item Explain what EEPROM is for and why it's unsuitable for data that changes every loop
    iteration.
  \item Read from and write to EEPROM at a given address.
  \item Explain why EEPROM's write sequence disables interrupts, and what could go wrong if it
    didn't.
  \item Explain why slow, blocking calls like \code{eeprom_write()} and \code{serial_print()}
    don't belong inside an ISR.
\end{itemize}

\needspace{6\baselineskip}
\subsection*{Questions to test yourself}
\begin{enumerate}
  \item Why does \code{eeprom_write_byte()} wait on the \code{EEPE} flag before starting a new
    write?
  \item What real AVR datasheet requirement forces \code{EEMPE} to be followed by \code{EEPE}
    within 4 clock cycles, and why does the example disable interrupts for that window?
  \item In \secref{c7:a:4}'s worked example, why does \code{ISR(PCINT0_vect)} only toggle LED1 and
    set a flag, leaving the actual \code{eeprom_write()} and UART report to \code{main()}?
\end{enumerate}

\needspace{6\baselineskip}
\subsection*{Looking ahead}
Chapter~\ref{c8:ch} is about the watchdog timer, another peripheral the chip has exactly one of,
and a hardware safety net against a program that hangs.
